\section{Introduction}
Speculative decoding changes how a language-model runtime executes work, not only how many target calls it makes. A draft proposes tokens, the target verifies several positions, and the runtime commits or rolls back state before continuing. The transformation provides useful acceleration only when the resulting execution satisfies its output requirement and the saved target work exceeds the added work~\cite{leviathan2023}. On quantized consumer systems, these conditions can fail separately.

For applications that must reproduce a selected greedy reference, equal committed token prefixes are an incomplete description of execution. Prefill and verification can use different numerical shapes, and previous calls can leave different numerical cache contents. A deviation at the next decision may depend on retained state, the current recipe, or their interaction. A repair chosen without separating these factors may reproduce one case and introduce another deviation. Even a repair that qualifies complete outputs can remove the performance advantage it was intended to preserve.

The same difficulty appears in the cost baseline. A repair can change target prefill or continuation before speculation is enabled. Comparing repaired speculation only with the original decoder can attribute this target-path improvement to speculation. Comparing only with a matched target-only path answers a different question and can hide the full deployment change. Both paid comparators are needed to interpret the result.

This paper examines these design decisions through controlled consumer-system studies. Early B570 and RTX interventions test whether output prediction establishes a useful repair. A matched RTX bridge tests whether a simpler draft-free execution explains the intervention output. Later natural-state crosses distinguish the state carried into a decision from the recipe used at that decision. A finite RTX attention-layout family then supplies a qualified repair, evaluated against ordinary and matched draft-free execution. Width tests, B570 costs, and historical M4 and 358H campaigns examine where qualification fails to produce a useful operating point.

The studies have different evidence roles. Selected-boundary interventions identify local effects; exposed panels guide repair selection; an input expansion checks a frozen recipe on project-unused inputs; paid sessions evaluate qualified outputs. We retain these distinctions rather than pooling all requests into a success rate. The four platforms are separate system cases, not four matched replications of one mechanism.

Our contribution connects execution-state consistency, repair scope, and paid system cost:
\begin{enumerate}
\item \textbf{Diagnostic evidence for separating state and current execution.} Natural donor-state crosses on CUDA and Vulkan reveal path effects, state effects, interactions, and selected-boundary nulls. Earlier prediction and simple-baseline results show why output reproduction alone cannot select a useful repair.
\item \textbf{A bounded repair-design case.} A query-layout constraint on target prefill and continuation is inclusion-minimal within a tested four-switch family on RTX. The frozen recipe qualifies complete reference responses on 24 project-unused MBPP inputs. This is a local graph transformation, not a portable determinism theorem.
\item \textbf{Qualified cost evidence and counterexamples.} RTX generation improves against both ordinary and matched draft-free execution. Much of its ordinary-reference request gain is already available without speculation. Wider fixed-capacity drafting and a B570 endpoint repair qualify while losing performance; M4 and 358H show limits on adaptation and draft-only optimization.
\end{enumerate}

The system-design lesson concerns three connected decisions: which execution state and path require control, which finite constraint preserves the registered output, and which baseline reveals operational value. Existing work establishes state consistency, numerical shape sensitivity, and verification overhead. We do not claim their discovery, a new general algorithm, or a universally optimal constraint. The additional evidence links these concerns within measured deployments and includes repairs that do not pay.

The paper defines the runtime and prior-work context, then the execution and measurement contracts. Its results follow diagnosis, repair design, and cost evaluation; the discussion extracts reusable design decisions while preserving deployment-specific limits.

